\section{Related Works}

\subsection{Speech Synthesis} 
With the rapid development of deep learning~\cite{DeT, zhao2025heterogeneous, yin2025progressive, chen2024sdpl,zhang2024inductive,tu2024smart, li2022long, zhang2025generating}, the FastSpeech series~\cite{fs2} first introduces a phoneme-level duration-based upsampling strategy and a controllable speech synthesis paradigm based on pitch and energy prediction.
% With recent advancements, speech synthesis evolve into two main paradigms: autoregressive and non-autoregressive models.
% Since VALL-E~\cite{zeroshotspeech} and SoundStream~\cite{soundstream} first introduced speech discretization and autoregressive generation paradigms into speech synthesis, research on autoregressive speech generation models and discrete speech codecs has continued to advance rapidly, such as SparkTTS~\cite{sparktts}, Llasa~\cite{llasa}, and CosyVoice series~\cite{cosyvoice1, cosyvoice2, cosyvoice3}.
Subsequently, many recent models, such as the StyleTTS~\cite{StyleTTS2} and NaturalSpeech series~\cite{naturalspeech3} achieve more natural speech synthesisby incorporating techniques such as diffusion models and adversarial training.
Meanwhile, speech synthesis models based on discrete speech codecs and autoregressive architectures have also emerged progressively, such as SparkTTS~\cite{sparktts}, Llasa~\cite{llasa}, and CosyVoice series~\cite{cosyvoice3}.
% Then, more recent methods such as StyleTTS~\cite{styletts, StyleTTS2} and NaturalSpeech~\cite{naturalspeech, naturalspeech2,naturalspeech3} families, there has been continuous progress in improving the naturalness and realism of generated speech, as well as the effectiveness of voice cloning.
Despite the progress, they cannot be directly applied to movie dubbing tasks because they lack the design of modeling duration and prosody from performance in the given movie clips.

\subsection{Visual Voice Cloning} 
Some previous V2C methods attempt to improve dubbing quality by pretraining the phoneme encoder~\cite{spk2dub} or decoupling acoustic modeling and prosody adaptation~\cite{ProDubber}, in response to the scarcity and noisiness of movie dubbing datasets caused by issues such as copyright constraints.
Another group of work explores techniques such as flow matching~\cite{emodubber} and contrastive learning~\cite{styledubber}, primarily aiming to enhance the performance of audiovisual alignment~\cite{HPMDubbing, zhao2025towards, cong2025flowdubber, li2025dubbing}.
However, their visual-based alignment methods struggle to generalize to unseen video scenarios, hindering the broader application.

% The V2C task requires generating high-quality dubbing speech how a text might be said, but in step with the lip movements portrayed by video character, and in vocal style exemplified by reference audio~\cite{V2C}.
% Some works focus primarily on improving the pronunciation clarity~\cite{HPMDubbing, styledubber, spk2dub}. 
% For example, SOTA dubbing method ProDubber~\cite{ProDubber} and Speaker2Dubber~\cite{spk2dub} propose a pre-training framework to learn clear pronunciation representation from a large-scale text-to-speech corpus~\cite{libritts}.  
% However, they over-rely on the TTS architecture and use an inaccurate duration predictor~\cite{ProDubber} to estimate the lip speaking time, without considering the intrinsic connection between visual movement and speech content. 
% Besides, StyleDubber~\cite{styledubber} attempts to use time stretching to balance articulation and lip-sync.
% Although the overall length can remain consistent, this does not fundamentally achieve the audio-visual alignment in fine-grained lip-sync. 
% In this work, we propose \dubber, a novel dubbing architecture that combines LLM-based semantic-aware learning with dual contrastive aligning to enhance the audio-visual sync and pronunciation, while improving acoustic quality by voice enhanced flow matching. 

\subsection{Dubbing with MLLMs} 
Multimodal large language models (MLLMs) have powerful and generalizable capabilities for multimodal content understanding.
The earliest works leveraging MLLMs for dubbing primarily utilized them to align the overall duration between the dubbing and the video for multilingual dubbing~\cite{dubwise}.
Recently, VoiceCraft-Dub~\cite{voicecraftdub} attempted to autoregressively generate dubbing by feeding video frames and script text into an LLM-decoder. DeepDubber~\cite{deepdubber} attempts to leverage a chain-of-thought prompting strategy to guide the MLLMs generating coarse-grained information—such as scene type, character age, and gender—from the video to assist alignment.
However, they overlook the capability of fine-grained MLLM-generated dubbing instructions in alignment and their potential in zero-shot dubbing.